\documentclass[11pt]{article}

% English submission draft derived from the approved Vietnamese review copy.
% It is self-contained and compiles with Overleaf's default pdfLaTeX engine.
\usepackage[T1]{fontenc}
\usepackage[utf8]{inputenc}
\usepackage[a4paper,margin=2.5cm]{geometry}
\usepackage{amsmath}
\usepackage{array}
\usepackage{booktabs}
\usepackage{longtable}
\usepackage{graphicx}
\usepackage{microtype}
\usepackage{xurl}

\setlength{\parindent}{1.2em}
\setlength{\parskip}{0.35em}
\raggedbottom
\newcolumntype{L}[1]{>{\raggedright\arraybackslash}p{#1}}

\newcommand{\datasetrepository}{Zenodo}
\newcommand{\datasetdoi}{\url{https://doi.org/10.5281/zenodo.22074838}}
\newcommand{\datasetlicense}{Creative Commons Attribution 4.0 International (CC BY 4.0)}
\newcommand{\coderepository}{\url{https://github.com/MinhTrannnnn/2.5D-}}
\newcommand{\figureplaceholder}[2]{%
  \fbox{\parbox[c][#1][c]{0.92\linewidth}{\centering\textbf{Figure placeholder}\\[0.6em]#2}}%
}

\title{A paired multi-difficulty 2.5D terrain dataset for ground-robot path planning}
\author{%
  Tran Quang Minh\\
  \small Posts and Telecommunications Institute of Technology, Hanoi, Vietnam\\
  \small Correspondence: \texttt{MinhTQ.B22VT351@stu.ptit.edu.vn}
}
\date{}

\begin{document}
\maketitle

\section*{Abstract}

We present a dataset of 5,010 procedurally generated 2.5D height-field maps for
ground-robot path-planning research. The data are organised as 1,670 paired
Easy--Medium--Hard triplets from five terrain families; the three members of
each triplet retain the same base realisation and ten Start--Goal coordinates.
Each map contains a $128\times128$ elevation grid at 0.25-m cell spacing,
together with support elevation, footprint slope, roughness, step height,
traversability cost, and footprint-aware collision layers. The release comprises 50,100
map--task records, 10,020 terrain and analysis images, 30,060 planner-output
images, and 90,180 trial records from a fixed protocol for six planners. A JSON
schema, deterministic seeds, SHA-256 checksums, validation reports, source code,
and a version-pinned Python environment support inspection, reproduction, and
reuse of the data. A pickle-free NPZ companion provides the 13 raster layers
for faster numerical loading, while JSON remains the canonical map record.

\section{Background \& Summary}

Ground robots operating outside structured indoor environments must account
for more than obstacles on a horizontal grid. Surface inclination, local
height discontinuities, and roughness can determine whether a robot pose or
transition is feasible, while elevation change can alter the cost of an
otherwise collision-free route.

Existing off-road resources cover several complementary research tasks. RUGD,
RELLIS-3D, WildScenes, and GOOSE provide real sensor observations and semantic
annotations for unstructured environments\cite{rugd,rellis3d,wildscenes,goose}.
ORFD focuses on off-road free-space detection, TartanDrive on learned vehicle
dynamics, and STONE on voxel-level 3D traversability prediction
\cite{orfd,tartandrive,stone}. Among navigation benchmarks, Verti-Bench evaluates
mobility over extremely rugged off-road terrain in high-fidelity multi-physics
simulation\cite{vertibench}, while MarsPlanBench and MoonPlanBench evaluate
global planning on occupancy grids derived from Martian and Lunar digital terrain
models\cite{planetarypathbench}. BARN is more directly concerned with navigation
than the perception datasets, but represents synthetic planar occupancy and
configuration-space environments rather than height-field terrain\cite{barn}.
These resources are not interchangeable: the perception and dynamics datasets
preserve real sensing and vehicle-response complexity, whereas BARN supplies a
controlled end-to-end navigation testbed. Tables~\ref{tab:related_resources}
and~\ref{tab:related_release} therefore separate dataset structure from released
planning products rather than attempting to rank dataset quality.

{\footnotesize
  \setlength{\tabcolsep}{3pt}
  \renewcommand{\arraystretch}{1.16}
  \begin{longtable}{L{0.12\linewidth} L{0.17\linewidth} L{0.20\linewidth}
                   L{0.26\linewidth} L{0.15\linewidth}}
  \caption{Structural comparison with related off-road datasets and navigation
  benchmarks. Scale denotes the primary released or annotated unit, so sensor
  frames and navigation maps are not treated as equivalent samples. ``n/a''
  means not applicable; ``n/r'' means not reported by the cited source.}
  \label{tab:related_resources}\\
    \toprule
    Resource & Origin and representation & Scale & Traversability or robot
    conditioning & Paired difficulty \\
    \midrule
    \endfirsthead
    \multicolumn{5}{c}{Table~\thetable{} (continued)}\\
    \toprule
    Resource & Origin and representation & Scale & Traversability or robot
    conditioning & Paired difficulty \\
    \midrule
    \endhead
    \midrule
    \multicolumn{5}{r}{Continued on next page}\\
    \endfoot
    \bottomrule
    \endlastfoot
    RELLIS-3D\cite{rellis3d} & Real RGB and LiDAR & 13,556 LiDAR scans; 6,235
    images (maps/tasks: n/a) & Semantic terrain labels; no footprint conditioning & No \\
    ORFD\cite{orfd} & Real RGB and LiDAR & 12,198 synchronised pairs
    (maps/tasks: n/a) & Three free-space classes; no explicit footprint model & No \\
    RUGD\cite{rugd} & Real RGB video & 7,453 annotated frames (maps/tasks: n/a) &
    Semantic terrain labels; no footprint conditioning & No \\
    WildScenes\cite{wildscenes} & Real images and LiDAR point clouds & 9,306
    annotated images; 12,148 point-cloud samples (maps/tasks: n/a) & Semantic
    labels; no footprint conditioning & No \\
    GOOSE\cite{goose} & Real RGB/NIR images and point clouds & 10,000 labelled
    pairs in the paper release (maps/tasks: n/a) & Semantic labels; no footprint
    conditioning & No \\
    TartanDrive\cite{tartandrive} & Real multimodal driving trajectories,
    including height maps & About 200,000 interactions (map/task counts n/r) &
    Vehicle-response data, not a traversability grid & No \\
    STONE\cite{stone} & Real cameras, LiDAR, radar, and 3D voxels & 43 sequences;
    50,878 frames (map/task counts n/r) & Geometry-aware 3D traversability labels & No \\
    BARN\cite{barn} & Synthetic 2D occupancy and configuration-space maps & 300
    maps; 300 predefined Start--Goal tasks & Robot-size-adaptable configuration
    space & Ordered independent environments, not paired \\
    Verti-Bench\cite{vertibench} & Simulated 2.5D elevation, semantic, and
    terrain-physics maps & 100 maps; 1,000 Start--Goal tasks & Nine vehicle types
    and multi-physics vehicle--terrain interaction & No; elevation levels are
    sampled independently \\
    MarsPlanBench/\newline MoonPlanBench\cite{planetarypathbench} & Mission-derived 2D occupancy
    grids from planetary DTMs/DEMs & More than 2,000 Mars maps; 36 Moon maps;
    task count n/r & Slope-threshold traversability; no rover footprint or
    dynamics & No; multiple slope thresholds are not paired-difficulty groups \\
    This release & Synthetic 2.5D height fields; JSON records and NPZ raster
    companions &
    5,010 maps; 16,700 unique paired tasks; 50,100 map--task records &
    Footprint-aware slope, roughness, step, cost, and collision layers & Yes:
    1,670 Easy--Medium--Hard triplets \\
  \end{longtable}
}

{\footnotesize
  \setlength{\tabcolsep}{4pt}
  \renewcommand{\arraystretch}{1.16}
  \begin{longtable}{L{0.18\linewidth} L{0.44\linewidth} L{0.29\linewidth}}
  \caption{Released planning products, fixed splits, and availability for the
  resources in Table~\ref{tab:related_resources}. ``Licence n/r'' means that
  the cited source did not report a licence; it does not imply unrestricted use.}
  \label{tab:related_release}\\
    \toprule
    Resource & Planner outputs and fixed splits & Data, code, and licence \\
    \midrule
    \endfirsthead
    \multicolumn{3}{c}{Table~\thetable{} (continued)}\\
    \toprule
    Resource & Planner outputs and fixed splits & Data, code, and licence \\
    \midrule
    \endhead
    \midrule
    \multicolumn{3}{r}{Continued on next page}\\
    \endfoot
    \bottomrule
    \endlastfoot
    RELLIS-3D\cite{rellis3d} & No fixed Start--Goal or planner outputs;
    split: yes & Data and code; CC BY-NC-SA 3.0 \\
    ORFD\cite{orfd} & Free-space benchmark; no route suite; split: n/r & Data
    and code; licence n/r \\
    RUGD\cite{rugd} & No route suite; segmentation split: yes & Data and splits;
    licence n/r \\
    WildScenes\cite{wildscenes} & No routes; segmentation split: yes;
    evaluation code & Data DOI and code; CC BY-NC-SA 4.0 \\
    GOOSE\cite{goose} & No route suite; train/validation/test split: yes & Data:
    CC BY-SA 4.0; code: MIT \\
    TartanDrive\cite{tartandrive} & No planner outputs; dynamics train/test
    split: yes & Data and code; CC BY 4.0 \\
    STONE\cite{stone} & Voxel-prediction benchmark; no fixed route suite;
    split: n/r & Paper/code public; dataset announced; CC BY-NC-ND 4.0 \\
    BARN\cite{barn} & Paths, difficulty metrics, and baseline results; no fixed
    train/test split & Data and code; licence n/r \\
    Verti-Bench\cite{vertibench} & Ten-system results over all tasks; no task
    used for training & Data and code; MIT for code, third-party components under
    their own licences \\
    MarsPlanBench/\newline MoonPlanBench\cite{planetarypathbench} & Classical- and learned-planner
    results; split: n/r & Public code and data; licence n/r \\
    This release & Fixed shared tasks, six-planner outputs, a fixed protocol,
    and group-level train/validation/test splits & Data: CC BY 4.0; code: MIT \\
  \end{longtable}
}

Against this context, the release provides a complementary, controlled
procedural design for grid-based ground-robot planning: large-scale terrain
morphology and navigation tasks are held fixed within each triplet while
geometric difficulty changes according to documented rules. The central unit
is an Easy--Medium--Hard triplet that shares a base terrain realisation and
Start--Goal tasks. Each map also includes footprint-aware analysis layers and
fixed-seed outputs from reference planners. The controlled pairing is the main
distinction from independently sampled or merely difficulty-ordered
environments; it permits within-base comparisons while holding terrain family,
seed-derived morphology, and task coordinates constant.

The dataset represents terrain as a 2.5D height field sampled on a grid: every
horizontal location $(x,y)$ has exactly one elevation value $z$, with no
missing cells. This representation supports geometry analysis at the robot-
footprint scale and direct use by grid-based or sampling-based planners. By
design, it cannot represent multiple surfaces at one horizontal coordinate,
such as bridges, overhangs, or caves. The five terrain families are
\texttt{smooth\_obstacles}, \texttt{rolling}, \texttt{mountain},
\texttt{rugged}, and \texttt{plateau}; they provide different combinations of
smooth mounds, rolling topography, ridges, locally rough regions, and
plateau-like structures.

Each base terrain realisation is generated independently from a seed and then
converted into Easy, Medium, and Hard members. The members retain the same
large-scale terrain structure and ten Start--Goal coordinates, but differ in
the relief-severity parameter and local surface detail. They are therefore paired observations rather
than independent maps. Difficulty-acceptance criteria constrain the navigable
fraction, the size and share of the largest connected component, the number of
connected components, and collision-boundary density. Construction also
requires the relief-severity parameter $q$ to increase and navigable fraction
to decrease from Easy to Hard.

The release contains 1,670 paired triplets and 5,010 map files. Each triplet
contributes ten shared tasks: three short, four medium, and three long tasks.
This gives 16,700 unique paired tasks and 50,100 map--task records. Within each
triplet, one shared long task is marked as canonical and appears in all three
maps; this task is used in the six-planner benchmark. Raster layers, tasks,
configurations, generation seeds, map-level metadata, visualisations, planner
outputs, integrity checklists, and a machine-readable benchmark protocol are
released together. This organisation
supports the broader goal of making research data findable and reusable through
clear metadata and documented processing procedures\cite{fair}.

The dataset can support controlled comparisons of navigation methods across
paired terrain difficulties, evaluation of terrain-aware heuristics and edge
costs, study of stochastic-planner stability under fixed budgets, and
development of data readers or learned traversability models. The benchmark
outputs are supplied as reference computations and technical checks of path
feasibility; they are not intended to establish a universally applicable
planner ranking.

\section{Methods}

\subsection{Computational design and spatial conventions}

The dataset was generated with Python 3.9.6, NumPy 2.0.2, and SciPy
1.13.1\cite{numpy,scipy}. The reference environment also pins Matplotlib 3.9.4
and Pillow 11.3.0 for image rendering. This release was generated, rendered,
and benchmarked on a MacBook Pro (model identifier Mac16,1) with an Apple M4
chip containing a 10-core CPU (four performance and six efficiency cores),
16 GB of unified memory, and macOS 26.5.2. Data generation used up to six worker
processes and the benchmark used three worker processes. This is the release-
generation configuration, not a minimum hardware requirement for reading the
data. Each map is a square raster of $128\times128$ cells at 0.25-m spacing,
corresponding to a nominal $32\times32$ m sampled region. Raster arrays are
stored as \texttt{array[y][x]}, whereas task endpoints are stored as
\texttt{[x,y]}. A grid cell is converted to metric coordinates by
\begin{equation}
  X=0.25x, \qquad Y=0.25y, \qquad Z=z(y,x),
\end{equation}
where $z$ is the raw surface-elevation layer in metres. Cell-centre coordinates
therefore range from 0 to 31.75 m on each horizontal axis.

Generation uses the master seed \texttt{7300}. For family index
$f\in\{0,\ldots,4\}$ and within-family sample index
$i\in\{1,\ldots,334\}$, the base seed is
\begin{equation}
  s_{f,i}=7300+1{,}000{,}003f+10{,}007(i-1).
\end{equation}
If a realisation does not satisfy all acceptance constraints, the program
makes up to ten deterministic attempts; attempt $a$ uses
$s_{f,i}+104{,}729a$. The accepted seed, task seed, configuration, and realised
parameters are stored in each map file.

\subsection{Base terrain families}

Every terrain family is a single-valued height field covering the complete
grid; blocked areas still contain terrain values and are not missing cells.
Correlated random fields are produced by applying Gaussian filters to normally
distributed noise and then normalising the variance. Analytic elliptical
features are subsequently placed with positions, axes, orientations, and
amplitudes derived from the seed. The same seed, configuration, source-code
version, and library versions reproduce the same surface. The five released
families combine multiscale correlated random fields with analytic geometric
elements.

Let $\mathcal{G}_{\sigma}$ denote an independent array of standard-normal
samples after a Gaussian filter of standard deviation $\sigma$ cells, followed
by subtraction of its sample mean and division by its sample standard
deviation. Filtering uses reflected boundaries. For \texttt{rolling},
\texttt{mountain}, and \texttt{rugged}, the stochastic background is
$\beta_f\mathcal{G}_{L/7}+0.42\beta_f\mathcal{G}_{L/18}$, where $L=128$ is the
grid width. Hill centres are sampled from
$U(-0.36E,0.36E)^2$, where $E=31.75$ m is the sampled horizontal extent. Each
hill is a rotated Gaussian ellipse with radii
$r_x\sim U(1.3,4.8)$ m and $r_y\sim U(1.2,4.4)$ m, orientation
$U(0,\pi)$, and amplitude $U(-0.55,1.0)A_{h,f}$. Ridges use radii
$U(4.0,9.0)$ m and $U(0.55,1.35)$ m and amplitude
$U(0.65,1.0)A_{r,f}$. Soft-edged mesas use radii $U(1.4,3.5)$ m and
$U(1.3,3.1)$ m, amplitude $U(1.3,3.2)$ m, and logistic edge softness
$U(0.07,0.16)$. Each roughness patch uses radii $U(1.5,4.2)$ m and
$U(1.4,3.8)$ m and multiplies a shared $\mathcal{G}_{0.8}$ field by a Gaussian
ellipse and $A_{p,f}$. Finally, independent tilts
$U(-0.018,0.018)$ m/m are applied along both horizontal axes. The
centres of ridges, mesas, and roughness patches follow the same distribution
as hill centres, and their orientations follow $U(0,\pi)$. The family-specific
constants are listed in Table~\ref{tab:family-parameters}.

\begin{table}[htbp]
  \centering
  \caption{Base-profile constants for the three multifeature terrain families.
  Counts are deterministic within a family; positions, orientations, radii,
  and amplitudes follow the distributions stated in the text. All amplitude
  constants are in metres.}
  \label{tab:family-parameters}
  \footnotesize
  \begin{tabular}{lccccc}
    \toprule
    Family & $\beta_f$ & \shortstack{Hills\\$(N,A_{h,f})$} &
    \shortstack{Ridges\\$(N,A_{r,f})$} & \shortstack{Mesas\\$N$} &
    \shortstack{Rough patches\\$(N,A_{p,f})$} \\
    \midrule
    \texttt{rolling}  & 0.24 & $(9,2.2)$  & $(1,2.4)$ & 1 & $(2,0.10)$ \\
    \texttt{mountain} & 0.38 & $(12,3.8)$ & $(3,4.2)$ & 1 & $(2,0.14)$ \\
    \texttt{rugged}   & 0.48 & $(13,2.8)$ & $(2,3.2)$ & 2 & $(6,0.26)$ \\
    \bottomrule
  \end{tabular}
\end{table}

The \texttt{smooth\_obstacles} background is
$0.055\mathcal{G}_{L/12}$. It contains three clustered domes placed at
offsets $-3.8$, 0, and 3.8 m along an axis with angle $U(0,2\pi)$, with
independent position jitter $U(-1,1)$ m on each coordinate, radii
$U(3.0,4.7)$ m and $U(3.0,4.8)$ m, and amplitudes
$U(2.6,3.9)$ m. Two isolated domes use both radii and amplitude from
$U(1.2,1.8)$ m. The cluster centre lies in
$[-0.22R,0.22R]\times[-0.08R,0.28R]$, and isolated-dome centres lie in
$[-0.62R,0.62R]\times[-0.68R,-0.36R]$, where $R=15.875$ m is the map
half-extent. Dome orientation is $U(-0.16,0.16)$ radians.

The \texttt{plateau} background is
$0.055\mathcal{G}_{L/12}+0.08\mathcal{G}_{L/16}$. The count is sampled
uniformly from \{3,4\}. Plateau centres are distributed around a phase
$U(0,2\pi)$
with angular jitter $U(-0.32,0.32)$ radians and radial distance
$U(0.28R,0.58R)$. Their two radii follow $U(2.3,4.8)$ m and
$U(2.0,4.3)$ m, orientation follows $U(0,\pi)$, height follows
$U(2.2,3.8)$ m, and logistic edge softness follows $U(0.075,0.13)$. All five
base-family surfaces are Gaussian-filtered with $\sigma=0.55$ cells and shifted
so that their minimum is zero.

\subsection{Paired difficulty construction}

For each base surface, a family-dependent local-detail field is constructed
with a generator seeded by the bitwise XOR of the realised seed and
\texttt{0x5F3759DF}. A
standard-normal field is Gaussian-filtered with correlation $\sigma_f$ and
divided by its sample standard deviation. Independently, a $U(0,1)$ field is
filtered with $\sigma=L/12$; values below its $\rho_f$ quantile are removed,
the remainder is normalised by its maximum and filtered once more with
$\sigma=1.2$. Their product is $\eta_f$. Table~\ref{tab:detail-parameters}
gives every family-specific value.

\begin{table}[htbp]
  \centering
  \caption{Parameters of the local-detail field and its severity-dependent
  amplitude. $s_f$ is the strength multiplier, $m_f$ the amplitude cap, and
  $\rho_f$ the envelope activation quantile.}
  \label{tab:detail-parameters}
  \small
  \begin{tabular}{lcccc}
    \toprule
    Family & $\sigma_f$ (cells) & $s_f$ & $m_f$ (m) & $\rho_f$ \\
    \midrule
    \texttt{smooth\_obstacles} & 2.8 & 0.24 & 0.045 & 0.58 \\
    \texttt{rolling}           & 2.3 & 0.38 & 0.070 & 0.52 \\
    \texttt{mountain}          & 1.7 & 0.58 & 0.105 & 0.44 \\
    \texttt{rugged}            & 1.2 & 0.78 & 0.145 & 0.34 \\
    \texttt{plateau}           & 2.8 & 0.28 & 0.055 & 0.58 \\
    \bottomrule
  \end{tabular}
\end{table}

A candidate difficulty surface is formed as
\begin{equation}
  z_d(x,y)=q z_0(x,y)+\alpha_f(q)\,\eta_f(x,y),
\end{equation}
where $z_0$ is the base realisation and the complete detail-amplitude rule is
\begin{equation}
  \alpha_f(q)=\min\!\left\{s_f\left[0.010q+
  0.060\max(0,q-0.75)^{1.24}\right],\ m_f\right\}.
\end{equation}
For severity interval $[q_{\min},q_{\max}]$, the exact 34 evaluated values are
not random; they are
\begin{equation}
  q_k=q_{\min}\left(\frac{q_{\max}}{q_{\min}}\right)^{k/33},
  \qquad k=0,\ldots,33,
\end{equation}
which is the definition used by \texttt{numpy.geomspace}. The general intervals are
0.18--1.05, 0.48--2.40, and 0.90--4.80 for Easy, Medium, and Hard,
respectively. \texttt{plateau} uses 0.18--1.25, 1.20--3.20, and 2.80--6.00
because a plateau top may remain a valid robot pose even when its surrounding
edge is impassable. If $q<0.62$, the candidate is additionally filtered with
$\sigma=1.15(0.62-q)$ cells. Its minimum is shifted to zero and elevation is
rounded to six decimal places before every derived layer and acceptance metric
is computed. The stored field \texttt{relief\_severity} is exactly the selected
multiplier $q$; it is not the realised max--min elevation range.

Candidate selection uses the footprint-aware collision grid described below.
Table~\ref{tab:difficulty-acceptance} gives all acceptance ranges and boundary
thresholds. Within every accepted triplet, $q$ strictly increases and the
navigable fraction decreases by more than 0.01 at each difficulty transition.
Candidate combinations are ranked by distance from the target navigable
fractions; the program examines up to the top 12 combinations and retains the
lowest-scoring combination among those for which all shared tasks can be
constructed.

\begin{table}[htbp]
  \centering
  \caption{Difficulty-acceptance criteria used during terrain generation. The
  first three rows apply to four terrain families; \texttt{plateau} uses
  separately calibrated ranges. The upper block gives the calibration bands;
  the lower block gives the corresponding structural constraints.}
  \label{tab:difficulty-acceptance}
  \small
  \setlength{\tabcolsep}{5pt}
  \begin{tabular}{llccc}
    \toprule
    Scope & Difficulty & $q_{\min}$--$q_{\max}$ &
    \shortstack{Navigable\\fraction} & Target \\
    \midrule
    General & Easy   & 0.18--1.05 & 0.72--0.94 & 0.84 \\
    General & Medium & 0.48--2.40 & 0.54--0.76 & 0.65 \\
    General & Hard   & 0.90--4.80 & 0.40--0.65 & 0.52 \\
    Plateau & Easy   & 0.18--1.25 & 0.80--0.90 & 0.84 \\
    Plateau & Medium & 1.20--3.20 & 0.73--0.82 & 0.78 \\
    Plateau & Hard   & 2.80--6.00 & 0.69--0.76 & 0.73 \\
    \bottomrule
  \end{tabular}

  \medskip
  \begin{tabular}{llcccc}
    \toprule
    Scope & Difficulty & $\lambda_{\min}$ &
    \shortstack{Maximum\\components} & $\psi_{\min}$ &
    $\delta_{B,\max}$ \\
    \midrule
    General & Easy   & $\geq0.65$ & 40 & 0.90 & 0.080 \\
    General & Medium & $\geq0.45$ & 60 & 0.84 & 0.085 \\
    General & Hard   & $\geq0.34$ & 75 & 0.80 & 0.090 \\
    Plateau & Easy   & $\geq0.70$ & 25 & 0.90 & 0.060 \\
    Plateau & Medium & $\geq0.62$ & 30 & 0.90 & 0.060 \\
    Plateau & Hard   & $\geq0.55$ & 35 & 0.88 & 0.060 \\
    \bottomrule
  \end{tabular}
\end{table}

For an $H\times W$ collision grid $G$, let
$F=\{(x,y):G(y,x)=0\}$ and let $L_F$ be the largest connected component under
the exact motion graph defined below. The released metrics are
\begin{equation}
  \phi=\frac{|F|}{HW},\qquad
  \lambda=\frac{|L_F|}{HW},\qquad
  \psi=\frac{|L_F|}{|F|}=\frac{\lambda}{\phi},
\end{equation}
where $\phi$ is navigable fraction, $\lambda$ is largest-component fraction,
and $\psi$ is largest-component share. Collision-boundary density counts only
horizontal and vertical cell adjacencies:
\begin{equation}
  \delta_B=\frac{
  \sum_{y=0}^{H-1}\sum_{x=0}^{W-2}[G_{y,x}\ne G_{y,x+1}]+
  \sum_{y=0}^{H-2}\sum_{x=0}^{W-1}[G_{y,x}\ne G_{y+1,x}]}
  {H(W-1)+(H-1)W}.
\end{equation}
Here brackets denote an indicator. An initial eight-neighbour binary-component
calculation gives $\widetilde\lambda$ for inexpensive screening; accepted
candidates are then recomputed using the full movement graph. For a difficulty
with navigable interval $[\ell,u]$, target $t$, and minimum component fraction
$\lambda_{\min}$, the candidate score is exactly
\begin{equation}
  S=|\phi-t|+4\{[\ell-\phi]_+ +[\phi-u]_+\}
  +3[\lambda_{\min}-\widetilde\lambda]_+,
\end{equation}
where $[x]_+=\max(0,x)$. Candidates are sorted by $S$, filtered by every limit
in Table~\ref{tab:difficulty-acceptance}, and the first six feasible candidates
per difficulty are retained. Before the full motion-graph calculation, the
inexpensive screen rejects candidates with
$\widetilde\lambda<0.95\lambda_{\min}$; the exact $\lambda_{\min}$ threshold is
still applied after connected components are recomputed. A triplet score is
$\sum_d S_d$.

\subsection{Robot-footprint-aware terrain analysis}

The reference robot is represented by a horizontal circular footprint of
radius 0.45 m. At 0.25-m cell spacing, the digital radius is
\begin{equation}
  R_c=\left\lceil\frac{0.45}{0.25}\right\rceil=2,
  \qquad
  \mathcal{D}=\{(i,j)\in\mathbb{Z}^2:i^2+j^2\leq R_c^2\}.
\end{equation}
Thus the mask contains exactly 13 offsets in a $5\times5$ window: the centre,
four axial offsets at radius one, four diagonal offsets $(\pm1,\pm1)$, and four
axial offsets at radius two. This upward rounding is conservative: the outer
axial samples lie 0.50 m from the centre although the nominal radius is 0.45 m.
The configuration was not derived from a specific physical robot. The four values---0.45-m radius, $28^\circ$ slope,
0.10-m roughness, and 0.65-m step height---were fixed before release generation
to define a reference robot and to produce both traversable and blocked regions
across all five terrain families. They are not universal mechanical limits for
ground robots. The nominal diameter is 0.90 m. The raw elevation layer is
released so that users can recompute all footprint-aware layers using geometry
and thresholds appropriate to another robot.

A least-squares plane is fitted over the circular window at every candidate
robot-centre location. The plane intercept is stored as
\texttt{support\_elevation}; the horizontal gradient magnitude is converted to
degrees and stored as \texttt{footprint\_slope\_degrees}. Roughness is the root
mean square of elevation residuals from the plane. Local step height is the
largest eight-neighbour elevation difference at central cells within the
digital mask. A point-gradient slope layer after light smoothing is also retained
for descriptive analysis, but it is not used to determine the traversability
grid or navigation cost.

Specifically, let $\mathcal{W}=\{(u_i,v_i,z_i)\}_{i=1}^{N}$ be samples in the
circular window, where $(u_i,v_i)$ are metric offsets from the window centre.
Because the circular mask is symmetric about its centre, the least-squares
plane $\widehat z_i=\bar z+a u_i+b v_i$ is computed as
\begin{equation}
  \bar z=\frac{1}{N}\sum_{i=1}^{N}z_i,\qquad
  a=\frac{\sum_i u_i z_i}{\sum_i u_i^2},\qquad
  b=\frac{\sum_i v_i z_i}{\sum_i v_i^2}.
\end{equation}
The support elevation at the centre is $\bar z$; footprint slope and RMS
roughness are, respectively,
\begin{equation}
  \theta=\frac{180}{\pi}\arctan\!\left(\sqrt{a^2+b^2}\right),\qquad
  r=\sqrt{\frac{1}{N}\sum_{i=1}^{N}
  \left(z_i-\bar z-a u_i-b v_i\right)^2}.
\end{equation}
In the implementation, local sums of $z$, $z^2$, $uz$, and $vz$ are evaluated
simultaneously at all cell centres using two-dimensional convolutions from
\texttt{scipy.ndimage}; point gradients, Gaussian filters, and local maxima are
also evaluated over complete NumPy arrays. This vectorisation avoids per-cell
Python loops when analysing 82,083,840 locations across 5,010 maps. Step height
uses eight shifted neighbour arrays and a footprint-scale maximum filter.

For slope $\theta$, roughness $r$, and step height $h$, the dimensionless
traversability cost is
\begin{equation}
  C_{\mathrm{trav}}=
  0.45\frac{\theta}{28^\circ}+
  0.25\frac{r}{0.10\ \mathrm{m}}+
  0.30\frac{h}{0.65\ \mathrm{m}}.
\end{equation}
Let $B(x,y)$ be a Boolean variable for a robot-centre pose. The exact
classification rule is
\begin{equation}
  B=(\theta>28^\circ)\ \lor\ (r>0.10\ \mathrm{m})\ \lor\
  (h>0.65\ \mathrm{m})\ \lor\ (C_{\mathrm{trav}}\geq1).
\end{equation}
A pose is traversable if and only if $B$ is false and the centre is outside the
boundary strip. Specifically, a centre is blocked when $x<R_c$, $x\geq W-R_c$,
$y<R_c$, or $y\geq H-R_c$. Plane-fit convolutions use reflected input at the
array boundary; shifted arrays and the footprint maximum used for step height
use nearest-edge extension. The explicit two-cell strip then overrides these
intermediate boundary values in the final collision grid. No additional
obstacle dilation is required to reproduce the reference configuration; adding
a safety margin or changing robot geometry defines a different navigation
configuration.

\subsection{Motion model and path objective}

Connectivity checks and reference routes use 8-connected movement between cell
centres without diagonal corner cutting. A transition is rejected if either
endpoint is blocked, a diagonal passes between blocked orthogonal neighbours,
or the absolute elevation difference divided by horizontal distance, computed
from \texttt{support\_elevation}, exceeds $\tan(28^\circ)$. For horizontal
length $d$, support-elevation change $\Delta z$, and slope weight $w=3$, the
shared edge cost is
\begin{equation}
  c=\sqrt{d^2+\Delta z^2}
  \left[1+w\left(\frac{|\Delta z|}{d}\right)^2\right].
\end{equation}
The path objective is the sum of $c$ over all transitions. The same movement-
validity rules are used for connectivity checks, all six planners, and path
validation. Cost $c$ is used to rank Start--Goal pairs, serves as the objective
minimised by Dijkstra and A*, weights PRM edges, is used to select parents and
rewire branches in RRT*, and is reported as the final cost of every successful
path. BFS minimises only the number of steps.
RRT-Connect expands its trees using Euclidean distance together with collision
and slope checks, so $c$ is computed only afterwards to report the path found.
For A*, the heuristic is the planar Euclidean distance to the Goal multiplied
by cell size; this is a lower bound on the remaining path cost under $c$.

\subsection{Start--Goal task construction}

Tasks are selected on the Hard member of each triplet and then copied to the
Easy and Medium members. Candidate endpoints lie in the largest valid connected
component, found by a queue-based breadth-first traversal of the 8-connected
graph using the documented movement rules. Candidate points are normally at
least 1.5 m, or six cells, from the boundary. For the ten-task protocol, all
candidate cells are used when there are at most 30; otherwise, exactly 30 are
sampled without replacement using the fixed seed. Dijkstra's
algorithm\cite{dijkstra} is run from each source cell
to compute weighted distances to other candidates; these distances rank
endpoint pairs under the shared path objective.

Under the ten-task protocol, three pairs are selected from distance quantiles
0.20--0.42, four from 0.45--0.72, and three from 0.75--0.98. They receive the
labels \texttt{short}, \texttt{medium}, and \texttt{long}, respectively. The
longest task is marked as the canonical visualisation and benchmark task. A
triplet is accepted only if every selected pair remains connected under the
documented movement model at all three difficulty levels.

For exact tie and shortage handling, endpoint pairs are stored as unordered
pairs and ranked first by weighted distance and then lexicographically by row
and column coordinates. For a ranked list of length $n$ and quantile interval
$[p_l,p_u]$, the slice begins at
$\min(n-1,\lfloor p_l(n-1)\rfloor)$ and ends before
$\min(n,\max(i_l+1,\lfloor p_u n\rfloor))$, where $i_l$ is the start index.
For a band with quota $m$, the algorithm takes $m$ equally spaced interior
positions between 0 and $|\mathrm{band}|-1$, rounds each with Python/NumPy's
nearest-even rule, and scans cyclically forward to the first pair not already used.
If a quantile slice contains fewer than $m$ candidates, the complete ranked
list replaces that slice. If the three bands still supply fewer than ten unique
pairs, unused pairs are taken from the complete list in decreasing-distance
order and labelled \texttt{long}. Generation of the current attempt fails if
ten distinct pairs still cannot be formed. If fewer than two margin-qualified
endpoint cells exist, the complete largest component is used; if it also has
fewer than two cells, the attempt fails.

\subsection{End-to-end matched-generation algorithm}

\noindent\textbf{Algorithm 1: Generate one paired terrain group.}
For family $f$ and within-family index $i$:
\begin{enumerate}
  \item Compute $s_{f,i}$ and, for attempt $a=0,\ldots,9$, set the realised
  seed to $s_{f,i}+104{,}729a$. Generate one base surface $z_0$ and one local
  detail field $\eta_f$ from this realised seed.
  \item For each difficulty, evaluate all 34 values $q_k$. For every value,
  construct $z_d$, apply the conditional smoothing and six-decimal rounding,
  compute all footprint-aware layers and metrics, calculate $S$, and retain the
  first six feasible candidates after sorting by $S$.
  \item Enumerate the Cartesian product of the three candidate lists. Discard
  combinations unless $q_{\mathrm{Easy}}<q_{\mathrm{Medium}}<q_{\mathrm{Hard}}$
  and $\phi_{d}-\phi_{d+1}>0.01$. Sort the rest by $\sum_d S_d$.
  \item For at most the first 12 combinations, construct ten tasks on the Hard
  member using the bitwise XOR of $(s_{f,i}+104{,}729a)$ and
  \texttt{0x6A09E667} as the task seed, following the deterministic rules
  above. Reject the combination if any task endpoints are
  disconnected in any of its three members.
  \item Return the first surviving combination and serialize its three maps
  with the same tasks. If none survives, continue with attempt $a+1$; report a
  generation error after ten unsuccessful attempts.
\end{enumerate}

\subsection{Reference planner benchmark}

The release contains a fixed-protocol benchmark for breadth-first search (BFS),
Dijkstra\cite{dijkstra}, A*\cite{astar}, a probabilistic roadmap (PRM)
implementation\cite{prm}, RRT-Connect\cite{rrtconnect}, and
RRT*\cite{rrtstar}. The benchmark evaluates the canonical long task on all
5,010 maps. This scope was fixed before execution to provide one consistently
defined long-route unit per map and to obtain complete coverage of all paired
difficulty members under the same computational protocol. It is a reference
validation benchmark, not an exhaustive evaluation of the ten-task suite. All
ten tasks remain published, and the released command-line program supports an
all-task run without regenerating terrain. Results reported here should not be
generalised to the other nine tasks. BFS is an unweighted reference that
minimises the number of transitions; Dijkstra and A* minimise the common
weighted objective. The deterministic planners run once per map. PRM,
RRT-Connect, and RRT* run with five fixed seeds; the system performs no
additional reruns beyond these five documented seeds.

PRM samples 500 free cells and considers the 12 nearest neighbours of each
node. RRT-Connect has a maximum budget of 5,000 iterations, a four-cell
extension step, and stops when the two trees connect or the budget is exhausted.
RRT* runs for all 6,000 iterations with a three-cell extension step, a seven-
cell rewiring radius, and a Goal-sampling probability of 0.08, then returns the
lowest-cost valid Goal connection it found. Edges from the sampling-based
planners are converted into symmetric 8-connected Bresenham cell sequences and
are accepted only when every rasterised transition passes the shared collision,
corner, and slope checks.

Every trial records success status, path validity, runtime, planar and three-
dimensional length, weighted cost, cumulative elevation gain, number of path
cells, and number of examined edges. The protocol produces 90,180 raw rows:
three deterministic and fifteen stochastic trial records per map. Summary
tables first group the five stochastic trials within each map--planner pair and
then aggregate across maps, avoiding treatment of five seeds on one terrain as
five independent terrain samples. Runtime was measured while the benchmark ran
in parallel with three worker processes and is stored with that execution
context.
These timings describe the completed three-worker throughput run only. They are
not used to rank planners or support inferential runtime comparisons, because
concurrent load was present and each deterministic planner has only one timing
observation per map. A timing study intended for performance comparison should
run planners sequentially, repeat each configuration, and report robust
summaries such as the median and interquartile range.

\subsection{Reproducibility and integrity-check preparation}

Data generation, image rendering, benchmarking, analysis, and validation are
implemented as separate command-line programs. The Conda environment and the
complete list of Python package versions accompany the source code. Generation
and benchmark jobs use deterministic configuration signatures and can resume
at triplet or completed-map granularity. The final validator checks the core
structure, \texttt{grid} and \texttt{elevation} arrays, recomputation of the
collision grid and \texttt{support\_elevation}, coordinate conventions, task
counts and connectivity, shared endpoints, and difficulty ordering. Path
validity is checked during benchmarking. Three SHA-256 lists cover all 5,010
JSON map files, 5,010 NPZ companions, and core release metadata. The complete
file inventory records relative paths,
content types, and byte sizes.

Figure~\ref{fig:terrain-triplets} illustrates how a base terrain realisation
and its Start--Goal task are retained across Easy, Medium, and Hard members of
each paired triplet.

\begin{figure}[htbp]
  \centering
  \IfFileExists{figures/terrain_triplet_001.png}{%
    \includegraphics[width=\linewidth,height=0.78\textheight,keepaspectratio]{figures/terrain_triplet_001.png}%
  }{%
    \figureplaceholder{75mm}{Upload
    \texttt{figures/terrain\_triplet\_001.png}.}%
  }
  \caption{\textbf{Procedurally generated terrains at three paired difficulty levels.}
  Columns show Easy, Medium, and Hard, and each row represents one terrain
  family. Members of a row share the base terrain realisation and Start--Goal
  coordinates. Grey regions are locations at which the reference robot centre
  is blocked. The layout combines 15 images rendered by the project pipeline;
  terrain data, blocked regions, and task coordinates were not altered.}
  \label{fig:terrain-triplets}
\end{figure}

\section{Data Records}

\subsection{Data release}

The complete dataset is archived on \datasetrepository{} at
\datasetdoi{}\cite{paired25ddata}. The corresponding source code is archived
at \coderepository{}. The complete
data release is approximately 32.6 GB and contains 50,136 inventoried files.
The source-code repository is intended for software and compact documentation,
not as a substitute for long-term archival of the complete dataset.

\subsection{Directory organisation}

The release root is \texttt{dataset\_5010\_v1/}. Map files and images are
partitioned by terrain family so that the three difficulty members of a paired
group remain close together while avoiding thousands of files in a single
directory. The principal structure is

\begin{verbatim}
dataset_5010_v1/
|-- maps/<family>/
|-- npz/<family>/
|-- images/
|   |-- terrain/<family>/
|   |-- analysis/<family>/
|   |-- paths/<family>/
|   `-- contact_sheets/
|-- metadata/
|   |-- dataset_summary.csv
|   |-- task_summary.csv
|   |-- splits/
|   |   |-- train.csv
|   |   |-- validation.csv
|   |   |-- test.csv
|   |   `-- split_protocol.json
|   |-- manifest.json
|   |-- validation_report.json
|   |-- pathfinding_benchmark_v1_trials.csv
|   |-- pathfinding_benchmark_v1_summary.csv
|   |-- pathfinding_benchmark_v1_protocol.json
|   |-- release_manifest.json
|   |-- release_filelist.csv
|   |-- npz_conversion.json
|   |-- npz_manifest.csv
|   |-- SHA256SUMS_CORE.txt
|   |-- SHA256SUMS_MAPS.txt
|   |-- SHA256SUMS_NPZ.txt
|   `-- benchmark_results/
|       |-- RESULTS.md
|       |-- benchmark_audit.json
|       |-- tables/
|       `-- figures/
|-- examples/load_map.py
|-- README.md
`-- SCHEMA.md
\end{verbatim}

Each family directory under \texttt{maps/} contains 1,002 JSON files: 334 base
samples multiplied by three difficulty levels. Filenames follow
\path{terrain_<family>_<index>_<difficulty>.json}; for example,
\path{terrain_mountain_001_hard.json}. The shared identifier
\path{terrain_mountain_001} links the Easy, Medium, and Hard members.

\subsection{Contents and structure of a JSON map file}

Each JSON file is a self-contained record of one terrain. It contains
identifiers, seeds, the raw elevation matrix, all derived terrain-analysis
layers, the final collision grid, ten task objects, coordinate and navigation
conventions, generation settings, difficulty-acceptance criteria, and map-level
metadata. Every raster layer is $128\times128$. Arrays use \texttt{[y][x]}
indexing, whereas task endpoints use \texttt{[x,y]}. In the binary collision
grid, 0 denotes a valid robot-centre cell and 1 denotes a blocked cell.

Table~\ref{tab:raster-fields} summarises the principal rasters. Raw elevation
is retained separately from \texttt{support\_elevation}: the former describes
the surface at each cell, whereas the latter is the intercept of the footprint-
scale support plane used for edge-slope and path-cost calculations. Cause masks
distinguish regions blocked by slope, roughness, step height, or combined
traversability cost.

\begin{table}[htbp]
  \centering
  \caption{Principal raster layers in each JSON map file.}
  \label{tab:raster-fields}
  \small
  \begin{tabular}{>{\ttfamily\raggedright\arraybackslash}p{0.28\linewidth}
                  >{\raggedright\arraybackslash}p{0.16\linewidth}
                  >{\raggedright\arraybackslash}p{0.48\linewidth}}
    \toprule
    \normalfont Field & Encoding or unit & Meaning \\
    \midrule
    elevation & m & Raw single-valued terrain elevation. \\
    grid & 0/1 & Final footprint-aware collision grid. \\
    point\_slope\_degrees & degrees & Gradient slope of the lightly smoothed surface. \\
    support\_elevation & m & Intercept of the least-squares support plane. \\
    footprint\_slope\_degrees & degrees & Inclination of the footprint-scale support plane. \\
    roughness & m RMS & Residual elevation variation about the support plane. \\
    step\_height & m & Largest adjacent elevation discontinuity within the footprint. \\
    traversability\_cost & dimensionless & Weighted sum of normalised slope, roughness, and step height. \\
    blocked\_by\_* & 0/1 & Cause masks for slope, roughness, step height, and combined cost. \\
    centre\_blocked & 0/1 & Union of the cause masks before the boundary constraint. \\
    \bottomrule
  \end{tabular}
\end{table}

Each task object contains an index, difficulty-specific and paired identifiers,
Start and Goal cells, distance class, reference weighted distance, reference
difficulty, and a canonical-task flag. For convenience, the canonical task is
also repeated in the top-level fields \texttt{start}, \texttt{goal},
\texttt{start\_z}, and \texttt{goal\_z}. The complete field dictionary,
coordinate conversions, units, and minimal parser are documented in
\texttt{SCHEMA.md}; \texttt{examples/load\_map.py} supplies a reader with no
external dependencies and basic structural checks.

\subsection{NPZ raster companions}

Each JSON map has a same-stem file under \texttt{npz/<family>/}. It contains
the 13 raster fields in Table~\ref{tab:raster-fields}, all with shape
$128\times128$: continuous layers use \texttt{float32}, and binary grids and
masks use \texttt{uint8}. The files contain no pickled objects and can be read
with \texttt{numpy.load(path, allow\_pickle=False)}. They are derived views for
efficient array access; JSON remains authoritative for identifiers, tasks,
generation settings, and metadata. Conversion details, paths, sizes, and
per-file hashes are recorded in \texttt{npz\_conversion.json},
\texttt{npz\_manifest.csv}, and \texttt{SHA256SUMS\_NPZ.txt}.

\subsection{Metadata and benchmark records}

The file \texttt{metadata/dataset\_summary.csv} is an index across all five
families, with one row per map containing identifiers, seeds, realised
generation parameters, data paths, and map-level metadata.
\texttt{metadata/task\_summary.csv} has one row per unique paired task and
therefore 16,700 rows; canonical-task definitions remain in the JSON map files.
\texttt{metadata/manifest.json} records the release scope, configuration, image
layout, and benchmark links.

The raw benchmark table contains 90,180 rows. Its columns identify the terrain,
paired group, task, planner, trial, deterministic seed, budget, success and
validity flags, runtime, path length, weighted cost, elevation gain, number of
path cells, and number of examined edges. The accompanying summary CSV groups
trials by planner, family, and difficulty. The protocol JSON stores movement
rules, planner budgets, numbers of stochastic trials, rendering scope, and
execution conditions so that numerical outputs are not separated from the
settings that produced them.

The \texttt{benchmark\_results/} directory contains map-level and aggregate
tables, four benchmark-summary figures, a reader-facing report, and a machine-
readable check report. These records support Technical Validation and later
reuse; they do not replace the raw trial CSV. Release provenance is supported
by \texttt{validation\_report.json}, three SHA-256 lists,
\texttt{release\_manifest.json}, and \texttt{release\_filelist.csv}.

\subsection{Fixed train, validation, and test indexes}

Three fixed indexes are provided in \texttt{metadata/splits/}. Assignment is
performed on \texttt{matched\_group\_id}, rather than on individual maps, so the
Easy, Medium, and Hard members of one base terrain and all ten shared tasks
always remain in the same split. The assignment is stratified by terrain
family. Within each of the five families, 234 groups are assigned to training,
50 to validation, and 50 to test, giving 1,170/250/250 groups and
3,510/750/750 maps overall.

The assignment is deterministic. Within each family, group identifiers are
ranked by the SHA-256 digest of the protocol name, seed 20260830, and
\texttt{matched\_group\_id}; the first 234, next 50, and final 50 groups are
allocated to train, validation, and test. The CSV files contain one row per
group and list all three map identifiers and relative paths. The accompanying
\texttt{split\_protocol.json} records the rule, seed, counts, file hashes, and
zero-leakage checks. These indexes are intended for learned methods; the
released reference-planner benchmark covers all maps and did not train a model.

\subsection{Image products}

For each of the 5,010 maps, \texttt{images/terrain/} contains one terrain image
and \texttt{images/analysis/} contains one traversability-analysis image,
giving 10,020 map images in total. Each map also has six planner-output images,
giving 30,060 files under \texttt{images/paths/}. For a stochastic planner, the
image uses the first successful trial in fixed seed order; if none of its five
trials finds a path within the budget, the image shows the resulting no-route
state. The corresponding status is retained in the benchmark CSV. Two contact
sheets provide quick views of terrain triplets and outputs on Hard maps. Images
are derived products for visual inspection; JSON maps and trial CSV files are
the authoritative numerical data.

Figure~\ref{fig:analysis-products} shows the content of analysis images rather
than only listing directory paths. All three examples are Hard maps with index
001 and were fixed in advance by family: \texttt{smooth\_obstacles},
\texttt{rugged}, and \texttt{plateau} illustrate three different
morphologies. Every image retains the pipeline's three-panel layout: height
field and contours; traversability cost with impassable regions; and the reason
each robot-centre pose is blocked.

\begin{figure}[htbp]
  \centering
  \textbf{(a) \texttt{smooth\_obstacles}}\par\smallskip
  \IfFileExists{figures/analysis_smooth_obstacles_hard.png}{%
    \includegraphics[width=0.98\linewidth]{figures/analysis_smooth_obstacles_hard.png}%
  }{\figureplaceholder{42mm}{Missing \texttt{smooth\_obstacles} analysis image.}}%
  \par\medskip
  \textbf{(b) Rugged}\par\smallskip
  \IfFileExists{figures/analysis_rugged_hard.png}{%
    \includegraphics[width=0.98\linewidth]{figures/analysis_rugged_hard.png}%
  }{\figureplaceholder{42mm}{Missing \texttt{rugged} analysis image.}}%
  \par\medskip
  \textbf{(c) Plateau}\par\smallskip
  \IfFileExists{figures/analysis_plateau_hard.png}{%
    \includegraphics[width=0.98\linewidth]{figures/analysis_plateau_hard.png}%
  }{\figureplaceholder{42mm}{Missing \texttt{plateau} analysis image.}}%
  \caption{\textbf{Terrain-analysis products for three illustrative families at
  Hard difficulty.} The panels in each row encode elevation, traversability
  cost, and footprint-aware blocking causes. Every image was rendered directly
  from its corresponding index-001 JSON map file; examples were not selected manually by
  visual appearance.}
  \label{fig:analysis-products}
\end{figure}

Figure~\ref{fig:planner-route-comparison} illustrates the six types of planner-
output images stored for the same map and Start--Goal coordinates. It describes
file contents and is not used to rank or compare planner performance.

\begin{figure}[htbp]
  \centering
  \begin{tabular}{cc}
    \includegraphics[width=0.46\textwidth]{figures/bfs_mountain_hard.png} &
    \includegraphics[width=0.46\textwidth]{figures/dijkstra_mountain_hard.png} \\
    \small (a) BFS & \small (b) Dijkstra \\[0.4em]
    \includegraphics[width=0.46\textwidth]{figures/a_star_mountain_hard.png} &
    \includegraphics[width=0.46\textwidth]{figures/prm_mountain_hard.png} \\
    \small (c) A* & \small (d) PRM \\[0.4em]
    \includegraphics[width=0.46\textwidth]{figures/rrt_connect_mountain_hard.png} &
    \includegraphics[width=0.46\textwidth]{figures/rrt_star_mountain_hard.png} \\
    \small (e) RRT-Connect & \small (f) RRT* \\
  \end{tabular}
  \caption{\textbf{Actual outputs from six planners on the same task.} All
  panels use \texttt{terrain\_mountain\_001\_hard} and the same Start--Goal
  coordinates. Displayed routes were produced by the benchmark system and do
  not cross blocked cells.}
  \label{fig:planner-route-comparison}
\end{figure}

\section{Technical Validation}

\subsection{Core structure and paired-design validation}

The release validator rereads all 5,010 JSON map files rather than checking only
summary metadata. Checks cover the core structure; shapes and value domains of
\texttt{grid} and \texttt{elevation}; recomputation of the collision grid
and \texttt{support\_elevation}; \texttt{[y][x]} and \texttt{[x,y]}
coordinate conventions; task counts; Start--Goal validity; connectivity under
the documented movement model; and consistency of task coordinates within each
triplet. The \texttt{validation\_report.json} file records 5,010/5,010 valid
maps and 50,100/50,100 valid map--task records, with no errors or warnings in
these checks.

All 5,010 NPZ companions were additionally checked for ZIP CRC integrity,
the exact 13-field set, NPY headers, $128\times128$ shapes, dtypes, and payload
lengths. The conversion report records no failed files.

A supplementary release check over metadata and benchmark results examines all
1,670 paired groups. All 1,670 groups have a strictly increasing
relief-severity parameter $q$ and
strictly decreasing navigable fraction in Easy--Medium--Hard order, with zero
violations of either condition. This check matters because the members of a
triplet are paired samples and must not be treated as three independent
realisations.

\subsection{Difficulty-stratification validation}

Figure~\ref{fig:difficulty-overview} extends the ordering check with descriptive
statistics over all 5,010 maps. Each point is the mean of 334 maps for one
family--difficulty combination, and each error bar is one sample standard
deviation across maps. From Easy to Hard, all five families show a decrease in
navigable area together with increases in $q$ and mean footprint slope. The
figure confirms that difficulty labels follow the documented calibration
criteria; it is not used to compare planner performance.

\begin{figure}[p]
  \centering
  \IfFileExists{figures/dataset_difficulty_overview.png}{%
    \includegraphics[width=\linewidth,height=0.68\textheight,keepaspectratio]{figures/dataset_difficulty_overview.png}%
  }{%
    \figureplaceholder{90mm}{Upload
    \texttt{figures/dataset\_difficulty\_overview.png}.}%
  }
  \caption{\textbf{Calibration of paired terrain difficulty across the complete
  dataset.} The four panels show navigable fraction, relief severity $q$, mean
  footprint slope, and connected-component count for five families and three
  difficulty levels. Points and error bars are the mean $\pm$ one sample
  standard deviation over 334 maps for each family--difficulty combination.}
  \label{fig:difficulty-overview}
\end{figure}

\subsection{Paired planning-difficulty validation}

To assess difficulty with quantities not used by the label-acceptance rules, we
used the canonical long task whose Start--Goal coordinates are shared by all
three members of a triplet. The statistical unit is one matched triplet, with
$n=334$ triplets per family. For stochastic planners, the five trials were first
aggregated into a map-level success rate. For each metric, the within-triplet
difference was then computed as $\Delta=\mathrm{Hard}-\mathrm{Easy}$.
Table~\ref{tab:paired-planning-effects} reports the paired mean difference and
the normal-approximation 95\% confidence interval
$\bar\Delta\pm1.96s_\Delta/\sqrt{n}$.

\begin{table}[htbp]
  \centering
  \caption{Change in planning difficulty from Easy to Hard. $\Delta C^*$ is the
  change in A* optimal weighted path cost. The other columns are changes in
  success rate in percentage points. Each cell is the paired mean [95\% CI]
  over 334 triplets.}
  \label{tab:paired-planning-effects}
  \footnotesize
  \setlength{\tabcolsep}{4pt}
  \renewcommand{\arraystretch}{1.15}
  \begin{tabular}{lcccc}
    \toprule
    Terrain family & $\Delta C^*$ & $\Delta$ PRM & $\Delta$ RRT-Connect & $\Delta$ RRT* \\
    \midrule
    \texttt{smooth\_obstacles} & \shortstack{7.45\\{[7.02, 7.87]}} &
      \shortstack{$-0.06$\\{$[-0.18, 0.06]$}} &
      \shortstack{0.00\\{[0.00, 0.00]}} & \shortstack{0.00\\{[0.00, 0.00]}} \\
    \texttt{rolling} & \shortstack{26.65\\{[25.05, 28.25]}} &
      \shortstack{$-39.64$\\{$[-43.78, -35.50]$}} &
      \shortstack{$-22.22$\\{$[-25.85, -18.58]$}} &
      \shortstack{$-27.90$\\{$[-31.80, -24.01]$}} \\
    \texttt{mountain} & \shortstack{26.04\\{[24.43, 27.65]}} &
      \shortstack{$-44.85$\\{$[-48.90, -40.80]$}} &
      \shortstack{$-25.57$\\{$[-29.16, -21.98]$}} &
      \shortstack{$-29.52$\\{$[-33.33, -25.71]$}} \\
    \texttt{rugged} & \shortstack{33.62\\{[31.87, 35.36]}} &
      \shortstack{$-54.31$\\{$[-58.45, -50.17]$}} &
      \shortstack{$-34.67$\\{$[-38.66, -30.69]$}} &
      \shortstack{$-42.46$\\{$[-46.75, -38.16]$}} \\
    \texttt{plateau} & \shortstack{4.37\\{[3.98, 4.75]}} &
      \shortstack{0.00\\{[0.00, 0.00]}} &
      \shortstack{0.00\\{[0.00, 0.00]}} & \shortstack{0.00\\{[0.00, 0.00]}} \\
    \bottomrule
  \end{tabular}
\end{table}

Optimal cost increases from Easy to Hard in all five families, with confidence
intervals excluding zero. Fixed-budget success decreases clearly for
\texttt{rolling}, \texttt{mountain}, and \texttt{rugged}, but remains saturated
near 100\% for \texttt{smooth\_obstacles} and \texttt{plateau}. Thus, the
difficulty stratification consistently increases route cost, whereas its effect
on success probability depends on terrain family and planner; the labels do not
imply that every planner must degrade on every family.

\subsection{Path and benchmark-output validation}

Every path returned by a planner is validated before its trial record is
written. The check covers endpoint cells, collisions, diagonal corner-cutting
rules, edge-slope limits, and 8-connected continuity after rasterisation. Among
90,180 trials, 83,998 returned valid paths and 6,182 did not find a path within
the fixed budget; no path marked successful was invalid. All 15,030 BFS,
Dijkstra, and A* trials succeeded on the 5,010 canonical tasks, confirming that
each benchmark task has at least one path on the reference grid. Consequently,
there is no task on which all six planners failed. No-path outcomes occur only
for the stochastic PRM, RRT-Connect, and RRT* implementations under finite
budgets; they do not indicate that a task is disconnected. These counts check
validity and confirm that the benchmark records both successful trials and
no-path outcomes under fixed budgets; they are not used to rank planners.

The release check reads the \texttt{success} and \texttt{path\_valid} flags in
all 90,180 trial rows and inventories 30,060 map--planner output units. No
\texttt{success=True} row has a \texttt{path\_valid} value other than
\texttt{True}, and no planner-output image is missing. In addition, A* and
Dijkstra have the same stored weighted cost to six decimal places on all 5,010
canonical tasks; the maximum stored absolute difference is 0. This is a cross-
check between two weighted graph-search implementations operating on the same
cost model.

Table~\ref{tab:validation-summary} summarises the scope and outcome of the main
technical checks.

\begin{table}[htbp]
  \centering
  \caption{Summary of technical checks on the release.}
  \label{tab:validation-summary}
  \small
  \begin{tabular}{L{0.43\linewidth}r r}
    \toprule
    Verification & Objects checked & Detected errors \\
    \midrule
    Core map-file structure & 5,010 maps & 0 \\
    Start--Goal task validity & 50,100 map--task records & 0 \\
    Easy--Medium--Hard ordering of $q$ & 1,670 triplets & 0 \\
    Navigable-fraction ordering & 1,670 triplets & 0 \\
    Cross-split \texttt{matched\_group\_id} leakage & 1,670 triplets & 0 \\
    Validity of paths reported as successful & 83,998 paths & 0 \\
    Completeness of planner-output images & 30,060 expected images & 0 \\
    A*--Dijkstra cost consistency & 5,010 canonical tasks & 0 \\
    \bottomrule
  \end{tabular}
\end{table}

For the planner-output-image row, 0 in the ``Detected errors'' column means that
all 30,060 expected images exist. Image existence does not imply that a planner
always found a path; success or no-path status must be read from the trial CSV.
The split audit additionally confirms that every one of the 1,670 groups occurs
exactly once across the three indexes and that each split has its prescribed
per-family count.

\subsection{Reproducibility}

Seeds and configuration signatures are retained in the data release; Python
and library versions are provided in the source-code repository. Regression
tests cover horizontal planes, recovery of support-plane slope, deterministic
task construction and stratification, symmetric edge rasterisation, corner-
cutting rules, cost-weight use, and A*/Dijkstra agreement. Three SHA-256 lists
allow changes to the 5,010 JSON maps, NPZ companions, and core metadata to be
detected after download.

\section{Usage Notes}

Raster layers are indexed as \texttt{[y][x]}, whereas Start--Goal cells are
stored as \texttt{[x,y]}; interchanging these conventions silently transposes a
map. The dependency-free reader in \texttt{examples/load\_map.py} demonstrates
the intended JSON access pattern. For numerical work, the same-stem NPZ file
loads the raster arrays directly; task and configuration metadata should still
be read from JSON. A \texttt{long} task is selected on the Hard member
from the 0.75--0.98 band of ranked weighted distances and then copied to the
paired Easy and Medium maps. The longest selected task is the canonical task.

The supplied collision and traversability layers describe a circular 0.45-m
reference footprint with the published geometric limits. Users evaluating a
different robot should recompute the footprint-aware layers from raw
\texttt{elevation}, update the footprint-safe border and movement checks, and
revalidate task connectivity. Dilating the released \texttt{grid} directly
would count the reference footprint twice unless an additional safety margin
is deliberately required.

For learned methods, users should follow the fixed split indexes and treat
\texttt{matched\_group\_id} as indivisible. For paired analyses, the matched
triplet is the statistical unit; the three difficulty members and repeated
stochastic-planner seeds are not independent terrain samples. The dataset
describes only static geometry through a single-valued 2.5D height field. It
does not model soil deformation, wheel slip, or vehicle dynamics; vegetation,
water, mud, or semantic surface properties; sensor noise or localisation
uncertainty; or bridges, overhangs, and multi-level surfaces at one horizontal
coordinate. The eight-connected motion model also omits vehicle heading,
turning radius, steering constraints, and platform-specific kinematics. All
data were generated at 0.25 m per cell, and behaviour under a change of spatial
resolution has not been validated. Because the terrains are synthetic,
similarity and transfer to real terrain have likewise not been quantified.
Reference robot thresholds and planner outputs should therefore not be
interpreted as universal operating limits or rankings.

\section{Data Availability}

The complete release, comprising canonical JSON map files, pickle-free NPZ
raster companions, Start--Goal tasks, fixed split indexes, derived images,
benchmark records, metadata, validation reports, a file inventory, and SHA-256
checksums, is archived at \datasetrepository{} and identified by
\datasetdoi{}\cite{paired25ddata}. The dataset is released under
\datasetlicense{}. The
source-code repository is not a substitute for long-term archival of the full
32.6-GB data release.

\section{Code Availability}

The Python code used to generate terrain, analyse traversability, select tasks,
run six planners, render images, validate records, and prepare the release is
available at \coderepository{}. The repository contains the version
corresponding to this review manuscript, the Conda environment,
\texttt{requirements-lock.txt}, regression tests, and reproduction
instructions. The source code is provided under the MIT License.

\begin{thebibliography}{99}

\bibitem{rugd}
Wigness, M. B., Eum, S., Rogers, J. G. III, Han, D. K. \& Kwon, H. A RUGD
dataset for autonomous navigation and visual perception in unstructured outdoor
environments. In \textit{2019 IEEE/RSJ International Conference on Intelligent
Robots and Systems (IROS)}, 5000--5007 (2019).
\url{https://doi.org/10.1109/IROS40897.2019.8968283}.

\bibitem{rellis3d}
Jiang, P., Osteen, P. R., Wigness, M. B. \& Saripalli, S. RELLIS-3D dataset:
data, benchmarks and analysis. In \textit{2021 IEEE International Conference on
Robotics and Automation (ICRA)}, 1110--1116 (2021).
\url{https://doi.org/10.1109/ICRA48506.2021.9561251}.

\bibitem{wildscenes}
Vidanapathirana, K. et al. WildScenes: a benchmark for 2D and 3D semantic
segmentation in large-scale natural environments. \textit{The International
Journal of Robotics Research} \textbf{44}, 532--549 (2025).
\url{https://doi.org/10.1177/02783649241278369}.

\bibitem{goose}
Mortimer, P. et al. The GOOSE dataset for perception in unstructured
environments. In \textit{2024 IEEE International Conference on Robotics and
Automation (ICRA)}, 14838--14844 (2024).
\url{https://doi.org/10.1109/ICRA57147.2024.10611298}.
Data and code: \url{https://github.com/FraunhoferIOSB/goose_dataset}.

\bibitem{orfd}
Min, C. et al. ORFD: a dataset and benchmark for off-road freespace detection.
In \textit{2022 International Conference on Robotics and Automation (ICRA)},
2532--2538 (2022).
\url{https://doi.org/10.1109/ICRA46639.2022.9812139}.

\bibitem{tartandrive}
Triest, S., Sivaprakasam, M., Wang, S. J., Wang, W., Johnson, A. M. \& Scherer,
S. TartanDrive: a large-scale dataset for learning off-road dynamics models. In
\textit{2022 International Conference on Robotics and Automation (ICRA)}
(2022). \url{https://doi.org/10.1109/ICRA46639.2022.9811648}.

\bibitem{stone}
Park, K. et al. STONE dataset: a scalable multi-modal surround-view 3D
traversability dataset for off-road robot navigation. \textit{arXiv preprint}
arXiv:2603.09175 (2026).
\url{https://doi.org/10.48550/arXiv.2603.09175}.

\bibitem{vertibench}
Xu, T. et al. Verti-Bench: A general and scalable off-road mobility benchmark
for vertically challenging terrain. In \textit{Robotics: Science and Systems}
(2025). \url{https://doi.org/10.15607/RSS.2025.XXI.138}.
Project, data and code: \url{https://people.cs.gmu.edu/~xiao/Research/Verti-Bench/}.

\bibitem{planetarypathbench}
Chanc\'{a}n, M., Banerjee, A. \& Nikolakopoulos, G. Planetary terrain datasets
and benchmarks for rover path planning. \textit{arXiv preprint}
arXiv:2512.21438 (2025).
\url{https://doi.org/10.48550/arXiv.2512.21438}.
Code and data: \url{https://github.com/mchancan/PlanetaryPathBench}.

\bibitem{barn}
Perille, D., Truong, A., Xiao, X. \& Stone, P. Benchmarking metric ground
navigation. In \textit{2020 IEEE International Symposium on Safety, Security,
and Rescue Robotics (SSRR)}, 116--121 (2020).
\url{https://doi.org/10.1109/SSRR50563.2020.9292572}.

\bibitem{fair}
Wilkinson, M. D. et al. The FAIR Guiding Principles for scientific data
management and stewardship. \textit{Sci. Data} \textbf{3}, 160018 (2016).
\url{https://doi.org/10.1038/sdata.2016.18}.

\bibitem{numpy}
Harris, C. R. et al. Array programming with NumPy. \textit{Nature}
\textbf{585}, 357--362 (2020).
\url{https://doi.org/10.1038/s41586-020-2649-2}.

\bibitem{scipy}
Virtanen, P. et al. SciPy 1.0: fundamental algorithms for scientific computing
in Python. \textit{Nature Methods} \textbf{17}, 261--272 (2020).
\url{https://doi.org/10.1038/s41592-019-0686-2}.

\bibitem{dijkstra}
Dijkstra, E. W. A note on two problems in connexion with graphs.
\textit{Numerische Mathematik} \textbf{1}, 269--271 (1959).
\url{https://doi.org/10.1007/BF01386390}.

\bibitem{astar}
Hart, P. E., Nilsson, N. J. \& Raphael, B. A formal basis for the heuristic
determination of minimum cost paths. \textit{IEEE Transactions on Systems
Science and Cybernetics} \textbf{4}, 100--107 (1968).
\url{https://doi.org/10.1109/TSSC.1968.300136}.

\bibitem{prm}
Kavraki, L. E., Svestka, P., Latombe, J.-C. \& Overmars, M. H. Probabilistic
roadmaps for path planning in high-dimensional configuration spaces.
\textit{IEEE Transactions on Robotics and Automation} \textbf{12}, 566--580
(1996). \url{https://doi.org/10.1109/70.508439}.

\bibitem{rrtconnect}
Kuffner, J. J. \& LaValle, S. M. RRT-Connect: An efficient approach to
single-query path planning. In \textit{Proceedings of the IEEE International
Conference on Robotics and Automation}, 995--1001 (2000).
\url{https://doi.org/10.1109/ROBOT.2000.844730}.

\bibitem{rrtstar}
Karaman, S. \& Frazzoli, E. Sampling-based algorithms for optimal motion
planning. \textit{The International Journal of Robotics Research} \textbf{30},
846--894 (2011). \url{https://doi.org/10.1177/0278364911406761}.

\bibitem{paired25ddata}
Tran, Q. M. A paired multi-difficulty 2.5D terrain dataset for ground-robot
path planning. \textit{Zenodo} (2026).
\url{https://doi.org/10.5281/zenodo.22074838}.

\end{thebibliography}

\section*{Author Contributions}

T.Q.M. conceived and designed the dataset; developed the software for terrain
generation, traversability analysis, task selection, planner benchmarking, and
validation; generated, validated, curated, and documented the data release;
performed the analyses and prepared the figures and tables; and drafted,
revised, and approved the manuscript.

\section*{Competing Interests}

The author declares no competing interests.

\section*{Funding}

This research received no external funding.

\end{document}
